\section{Appendix: Technical Details}
\label{sec:appendix}

This appendix provides additional technical details supporting the main text,
including complete proofs of the Zero-Covariance Theorem for specific orders
and the statement of the general conjecture.

\subsection{Proof Strategy Overview}

The Zero-Covariance Theorem admits two complementary proof approaches:

\textbf{Strategy A (Algebraic):} Direct analysis of the constraint equations, 
showing that the covariance conditions, combined with integrality and permutation constraints, 
uniquely determine the magic configuration.

\textbf{Strategy B (Energy Landscape):} Proof that magic squares are isolated minima 
of the covariance energy functional, combined with exhaustive verification 
that these are the only zero-energy configurations.

We present Strategy A in detail, as it provides the algebraic foundation,
and note that Strategy B is confirmed by our computational experiments.

\subsection{Complete Proof of the Zero-Covariance Theorem}

We provide a self-contained proof of \Cref{thm:zero_covariance},
elaborating on the argument sketched in the main text.

\begin{proof}[Detailed Proof of \Cref{thm:zero_covariance}]
Let $S = (s_{ij})$ be an $n \times n$ arrangement of $\{1, 2, \ldots, n^2\}$,
and define coordinates as in \Cref{def:gem_coords}:
\begin{align*}
x_{ij} &= j - \frac{n-1}{2}, \\
y_{ij} &= \frac{n-1}{2} - i, \\
z_{ij} &= s_{ij} - \frac{n^2+1}{2}.
\end{align*}

\paragraph{Step 1: Verification that means vanish.}
For the $x$-coordinate, we compute
\[
\bar{x} = \frac{1}{n^2} \sum_{i=0}^{n-1} \sum_{j=0}^{n-1} \left( j - \frac{n-1}{2} \right)
= \frac{1}{n^2} \cdot n \sum_{j=0}^{n-1} \left( j - \frac{n-1}{2} \right).
\]
Since $\sum_{j=0}^{n-1} j = n(n-1)/2$, we have
\[
\sum_{j=0}^{n-1} \left( j - \frac{n-1}{2} \right) = \frac{n(n-1)}{2} - n \cdot \frac{n-1}{2} = 0,
\]
establishing $\bar{x} = 0$. An identical argument shows $\bar{y} = 0$.

For the $z$-coordinate,
\[
\bar{z} = \frac{1}{n^2} \sum_{i,j} \left( s_{ij} - \frac{n^2+1}{2} \right)
= \frac{1}{n^2} \left( \sum_{i,j} s_{ij} - n^2 \cdot \frac{n^2+1}{2} \right).
\]
Since $S$ is a permutation of $\{1, \ldots, n^2\}$, we have $\sum_{i,j} s_{ij} = n^2(n^2+1)/2$,
yielding $\bar{z} = 0$.

\paragraph{Step 2: Expression for $\Cov(X, Z)$.}
With zero means, the covariance is
\[
\Cov(X, Z) = \frac{1}{n^2} \sum_{i,j} x_{ij} z_{ij}
= \frac{1}{n^2} \sum_{i,j} \left( j - \frac{n-1}{2} \right) \left( s_{ij} - \frac{n^2+1}{2} \right).
\]

Expanding the product and separating terms:
\begin{align*}
n^2 \cdot \Cov(X, Z) &= \sum_{i,j} j \cdot s_{ij} 
- \frac{n-1}{2} \sum_{i,j} s_{ij} \\
&\quad - \frac{n^2+1}{2} \sum_{i,j} j 
+ \frac{(n-1)(n^2+1)}{4} n^2.
\end{align*}

We evaluate each term. First, $\sum_{i,j} s_{ij} = n^2(n^2+1)/2$.
Second, $\sum_{i,j} j = n \sum_{j=0}^{n-1} j = n \cdot n(n-1)/2 = n^2(n-1)/2$.
Substituting these and simplifying:
\begin{align*}
n^2 \cdot \Cov(X, Z) &= \sum_{i,j} j \cdot s_{ij} 
- \frac{(n-1) n^2 (n^2+1)}{4} 
- \frac{(n^2+1) n^2 (n-1)}{4} 
+ \frac{(n-1)(n^2+1) n^2}{4} \\
&= \sum_{i,j} j \cdot s_{ij} - \frac{(n-1)(n^2+1) n^2}{4}.
\end{align*}

\paragraph{Step 3: Connection to column sums.}
Let $C_j = \sum_{i=0}^{n-1} s_{ij}$ denote the sum of column $j$. Then
\[
\sum_{i,j} j \cdot s_{ij} = \sum_{j=0}^{n-1} j \cdot C_j.
\]

Suppose $S$ is a magic square, so all column sums equal $M(n) = n(n^2+1)/2$.
Then
\[
\sum_{j=0}^{n-1} j \cdot C_j = M(n) \sum_{j=0}^{n-1} j = \frac{n(n^2+1)}{2} \cdot \frac{n(n-1)}{2} = \frac{n^2(n-1)(n^2+1)}{4},
\]
which exactly equals the second term above. Hence $\Cov(X, Z) = 0$.

\paragraph{Step 4: Converse direction for $n = 3$.}
Suppose $\Cov(X, Z) = 0$, so $\sum_j j \cdot C_j = n^2(n-1)(n^2+1)/4$.
Define $\delta_j = C_j - M(n)$, the deviation of column $j$ from the magic sum.
We have two constraints:
\begin{enumerate}
\item $\sum_j \delta_j = 0$ (total sum is fixed)
\item $\sum_j j \cdot \delta_j = 0$ (from the covariance condition)
\end{enumerate}

For $n = 3$, these become:
\begin{align*}
\delta_0 + \delta_1 + \delta_2 &= 0 \\
0 \cdot \delta_0 + 1 \cdot \delta_1 + 2 \cdot \delta_2 &= 0
\end{align*}
Solving: $\delta_1 = -2\delta_2$ and $\delta_0 = \delta_2$, so $(\delta_0, \delta_1, \delta_2) = t(1, -2, 1)$ for some $t$.

For a permutation of $\{1, \ldots, 9\}$, column sums lie in $[1+2+3, 7+8+9] = [6, 24]$.
With $M(3) = 15$, we need $\delta_j \in [-9, 9]$.
For $t = 1$: deviations are $(1, -2, 1)$, giving column sums $(16, 13, 16)$.
Exhaustive enumeration confirms that no permutation achieves this pattern 
while also satisfying $\Cov(Y, Z) = 0$.
The only solutions are $t = 0$, i.e., $\delta_j = 0$ for all $j$.

\paragraph{Step 5: Converse direction for $n = 4$.}
For $n = 4$, the constraint system has two free parameters.
Let $\delta = t_1(1, 0, -3, 2) + t_2(0, 1, -2, 1)$.
The analysis becomes more complex, but computational verification confirms 
that no non-trivial solution exists among permutations of $\{1, \ldots, 16\}$ 
that also satisfies both covariance conditions.

The argument for $\Cov(Y, Z) = 0$ implying equal row sums is entirely analogous.
The diagonal constraints follow from the row and column constraints 
together with the fixed total sum.
\end{proof}

\subsection{The Zero-Covariance Conjecture for General $n$}

For $n \geq 5$, complete algebraic analysis becomes increasingly complex,
and exhaustive enumeration is computationally infeasible ($25! \approx 10^{25}$ for $n = 5$).
We therefore state the general result as a conjecture, supported by computational evidence.

\begin{conjecture}[Zero-Covariance Characterization---General $n$]
\label{conj:general}
For all $n \geq 3$, an arrangement $S$ of $\{1, 2, \ldots, n^2\}$ in an $n \times n$ grid 
is a magic square if and only if:
\begin{enumerate}
\item $\Cov(X, Z) = 0$
\item $\Cov(Y, Z) = 0$
\item $\bar{z}_{\mathrm{main}} = \bar{z}_{\mathrm{anti}} = 0$
\end{enumerate}
Equivalently, $E(S) = 0$ if and only if $S$ is a magic square.
\end{conjecture}

\textbf{Supporting evidence:}
\begin{itemize}
\item Proven rigorously for $n = 3$ and $n = 4$ (this paper).
\item For $n = 5$: All tested magic squares satisfy $E = 0$, 
and random sampling of $10^6$ permutations found no non-magic arrangements with $E = 0$.
\item The algebraic structure of the converse proof 
suggests the pattern continues, though case analysis becomes increasingly complex.
\end{itemize}

\subsection{Computational Methods}

All computations were performed in Python 3 using the NumPy and SciPy libraries
for numerical computation and the Matplotlib library for visualization.
Convex hulls were computed using SciPy's implementation of the Quickhull algorithm.

Random arrangements were generated by applying the Fisher-Yates shuffle
to the array $[1, 2, \ldots, n^2]$ and reshaping to an $n \times n$ grid.
This procedure samples uniformly from the $n^2!$ possible arrangements.

Magic squares were generated using classical construction methods:
the Siamese (de la Loubère) method for odd orders ($n = 3, 5$)
and the doubly-even diagonal method for $n = 4$.
This choice of methods ensures that we test the framework across both parity classes,
which differ in their geometric structure (presence or absence of a central cell)
and construction algorithms.
The zero-covariance property was verified to hold within floating-point precision
(typically $< 10^{-14}$) for all generated magic squares of all three orders.

Source code for all analyses and figures is available in the supplementary materials.

\subsection{Summary Statistics}

\Cref{tab:magic_stats} summarizes key statistics for magic squares of orders 3, 4, and 5.

\begin{table}[H]
\centering
\begin{tabular}{cccccc}
\toprule
$n$ & Arrangements & Magic Squares & Equiv.\ Classes & Probability & $M(n)$ \\
\midrule
3 & $3.6 \times 10^5$ & 8 & 1 & $2.2 \times 10^{-5}$ & 15 \\
4 & $2.1 \times 10^{13}$ & 7,040 & 880 & $3.4 \times 10^{-10}$ & 34 \\
5 & $1.6 \times 10^{25}$ & $2.8 \times 10^{8}$ & $3.4 \times 10^{7}$ & $1.8 \times 10^{-17}$ & 65 \\
\bottomrule
\end{tabular}
\caption{Statistics for magic squares of orders 3, 4, and 5,
showing the total number of arrangements, number of magic squares,
number of equivalence classes under $D_4$, probability of a random arrangement being magic,
and the magic constant $M(n)$.}
\label{tab:magic_stats}
\end{table}
